\documentclass[11pt,letterpaper]{article}
\usepackage[T1]{fontenc}
\usepackage{lmodern}
\usepackage[margin=1in]{geometry}
\usepackage{amsmath,amssymb,amsthm,mathtools}
\usepackage{microtype}
\usepackage{booktabs,longtable,array}
\usepackage{xcolor}
\usepackage{enumitem}
\usepackage{fancyhdr}
\usepackage{xurl}
\usepackage{hyperref}
\usepackage{listings}
\definecolor{ink}{RGB}{25,48,73}
\definecolor{linkblue}{RGB}{20,80,125}
\hypersetup{colorlinks=true,linkcolor=linkblue,urlcolor=linkblue,citecolor=linkblue,
 pdftitle={Riemann's functional equation in intuitionistic primitive recursive arithmetic},
 pdfauthor={ChatGPT}}
\setlength{\headheight}{14pt}
\pagestyle{fancy}
\fancyhf{}
\fancyhead[L]{\small Riemann's functional equation in iPRA}
\fancyhead[R]{\small Quantitative constructive proof}
\fancyfoot[C]{\thepage}
\setlength{\parskip}{0.35em}
\setlength{\parindent}{0pt}
\setlist{itemsep=0.25em,topsep=0.4em}
\allowdisplaybreaks[2]
\numberwithin{equation}{section}
\newtheorem{theorem}{Theorem}[section]
\newtheorem{lemma}[theorem]{Lemma}
\newtheorem{proposition}[theorem]{Proposition}
\newtheorem{corollary}[theorem]{Corollary}
\theoremstyle{definition}
\newtheorem{definition}[theorem]{Definition}
\newtheorem{remark}[theorem]{Remark}
\newcommand{\ipra}{\mathsf{iPRA}}
\newcommand{\PR}{\mathsf{PR}}
\newcommand{\Q}{\mathbb Q}
\newcommand{\N}{\mathbb N}
\newcommand{\Z}{\mathbb Z}
\newcommand{\C}{\mathbb C}
\newcommand{\R}{\mathbb R}
\newcommand{\ii}{\mathrm i}
\newcommand{\e}{\mathrm e}
\newcommand{\normone}[1]{\lVert #1\rVert_1}
\newcommand{\ps}{\psi}
\newcommand{\gamrec}{R}
\newcommand{\Jh}{\widehat J}
\newcommand{\lceilog}{\lambda}
\newcommand{\code}[1]{\texttt{\detokenize{#1}}}
\lstset{basicstyle=\small\ttfamily,breaklines=true,columns=fullflexible,
 keepspaces=true,frame=single,rulecolor=\color{black!20},showstringspaces=false}

\begin{document}
\hypersetup{pageanchor=false}
\begin{titlepage}
\vspace*{0.65in}
{\color{ink}\Huge\bfseries Riemann's functional equation\par}
\vspace{0.3in}
{\LARGE in intuitionistic primitive recursive arithmetic\par}
\vspace{0.4in}
{\large A complete quantitative human proof and a Lean formalization plan\par}
\vspace{0.35in}
{\large Constructive proof report\par}
{\large 15 November 2026\par}
\vspace{0.5in}
\begin{abstract}
We construct the meromorphic Riemann zeta function by explicit primitive
recursive approximation procedures and prove
\[
 \zeta(s)=2^s\pi^{s-1}\sin(\pi s/2)\Gamma(1-s)\zeta(1-s).
\]
The primary arithmetic statement is an everywhere-defined,
denominator-cleared identity between specified complex rational
approximants at every requested precision. The usual equation is its
meromorphic interpretation.

The proof uses a quantitative finite Fej\'er calculation to establish the
Gaussian theta transformation, finite Euler products to establish the
reciprocal-gamma identities, and a rapidly convergent regularized Mellin
integral. Gaussian normalization follows from the theta calculation at
\(t=1\). All convergence, quadrature, and division steps have explicit
numerical controls. The iPRA interpretation uses symbolic numerical
parameters and quantifier-free induction on finite arithmetic invariants.

The mathematical proof is supplied here. The Lean part is a concrete
formalization blueprint with small interface drafts; it is not a completed
or compiled Lean formalization of the analytic argument.
\end{abstract}
\vfill
{\small Repository snapshot:
\code{VladimirReshetnikov/ProveIt} at
\code{c5093c20b5bffb0333fd42915b7bc858d5aa45cc}.\\
The report is self-contained in its analytic argument. Existing
ConciseIPRA results supply the intended object-theory interface.}
\end{titlepage}
\clearpage
\pagenumbering{roman}
\hypersetup{pageanchor=true}
\begingroup
\small
\setlength{\parskip}{0pt}
\tableofcontents
\endgroup
\newpage
\pagenumbering{arabic}

\section{The result and its precise arithmetic meaning}

The familiar functional equation is an identity of meromorphic functions,
not an instruction to multiply separately undefined values at their poles.
We prove it through an identity of entire numerators. This also gives the
appropriate first-order arithmetic formulation.

Put
\begin{align}
 \ps(t)&=\sum_{n=1}^{\infty}\e^{-\pi n^2t}, &&t>0,\label{eq:psi}\\
 J(s)&=\int_1^\infty \ps(t)t^{s/2-1}\,dt,\label{eq:J}\\
 X(s)&=\frac12+\frac{s(s-1)}2\bigl(J(s)+J(1-s)\bigr),\label{eq:X}\\
 \gamrec_N(z)&=zN^{-z}\prod_{n=1}^N(1+z/n),\qquad
 \gamrec(z)=\lim_{N\to\infty}\gamrec_N(z),\label{eq:R}\\
 Z(s)&=\pi^{s/2}\gamrec(1+s/2)X(s).\label{eq:Z}
\end{align}
Powers of positive real numbers always mean \(a^z=\exp(z\log a)\).
There is no choice of a complex logarithm. The limits and integrals in
these displays will be given specified rational approximation algorithms.

\begin{theorem}[Constructive functional equation]\label{thm:main}
The functions \(J,X,\gamrec,Z\) in
\eqref{eq:J}--\eqref{eq:Z} have uniformly primitive recursive rational
approximants on each explicitly bounded rational box. They are entire in
the ordinary analytic interpretation, and
\begin{align}
 X(1-s)&=X(s),\label{eq:xiFE}\\
 s\gamrec(1-s)Z(s)
 &=(1-s)\,2^s\pi^{s-1}\sin(\pi s/2)Z(1-s) \qquad(s\in\C).
 \label{eq:entireFE}
\end{align}
The meromorphic function represented by
\[
                    \zeta(s)=\frac{Z(s)}{s-1}
                    \tag{\(\ast\)}
\]
agrees with \(\sum_{n\ge1}n^{-s}\) for \(\Re s>1\), has its only pole at
\(s=1\), with residue \(1\), and satisfies the requested equation
\begin{equation}
 \boxed{\zeta(s)=2^s\pi^{s-1}\sin(\pi s/2)\Gamma(1-s)\zeta(1-s)}
 \label{eq:standardFE}
\end{equation}
as a meromorphic identity. Here \(\Gamma\) is represented meromorphically
by \((1,\gamrec)\), and agrees with Euler's integral on \(\Re z>0\).

With the interpretation in Section~\ref{sec:arithmetic}, the
finite-precision statements of this theorem have proofs in \(\ipra\).
They also specialize to every fixed primitive recursive complex Cauchy
name supplied with the stated modulus and a numerical bound.
\end{theorem}

The notation \(X\) denotes the conventional Riemann \(\xi\)-function.
We use \(Z\) for its entire zeta numerator, so \(Z\) is not the Dirichlet
series itself. Formula~\eqref{eq:standardFE} and its equivalence to
\(\xi(s)=\xi(1-s)\) are the standard normalizations
\cite{DLMF25}. The contribution of this report is the explicit weak-arithmetic
proof and its finite interpretation.

\subsection{What is quantified inside iPRA}

First-order \(\ipra\) has natural numbers, primitive recursive function
symbols, their defining equations, intuitionistic first-order logic,
successor axioms, and quantifier-free induction. It has no sort of all
real numbers. Accordingly our principal internal theorem has the form
\begin{equation}
 \ipra\vdash
 \forall c\,\forall B\,\forall k\;
                \mathsf{Check}(c,B,k)=1.                     \label{eq:checktarget}
\end{equation}
Here \(c\) codes a rational complex input, \(B\) is a box bound, \(k\) is
the requested accuracy, and \(\mathsf{Check}\) is one fixed primitive
recursive function. On valid inputs it compares rational approximants
to the two sides of~\eqref{eq:entireFE}; on malformed or out-of-box
inputs it returns \(1\). Section~\ref{sec:arithmetic} specifies it.
Thus no undefined value or undecidable real-zero test occurs.

An equality involving real or complex names abbreviates the family of
finite rational estimates at every precision. We prove those estimates
with their numerical parameters free and then universally generalize.
We do not induct on the quantified assertion that a name is Cauchy or
that two names are equal.

For a fixed PR name \(a:\N\to\Q(\ii)\), the assertion
\[
 \normone{a(p)-a(q)}\le 2^{-p}+2^{-q}
\]
is a regularity hypothesis; operations on that name are compiled from
its fixed PR definition. No universal PR evaluator for arbitrary PR
program codes is assumed. Producing separate numerical proof
certificates also would not by itself establish~\eqref{eq:checktarget};
our induction arguments retain the precision as a symbolic parameter.

\section{The finite analytic toolkit}\label{sec:toolkit}

We give the elementary infrastructure explicitly, since its quantitative
content is essential here. The notation \(|z|\) is the usual complex
modulus in the human proof. The rational implementation uses
\(\normone{x+\ii y}=|x|+|y|\), with
\[
 |z|\le\normone z\le2|z|,\qquad
 \normone{zw}\le\normone z\,\normone w.
\]
Every final inequality can therefore be decided using rational
arithmetic without square roots.

\subsection{Elementary functions and constants}

Define \(\exp z\) by its factorial series and \(\sin,\cos\) by the
corresponding series. On \(\normone z\le S\), for an integer \(S\ge1\),
the exponential tail after degree \(m\ge2S\) is at most
\begin{equation}
              2^{\,4S-m}.                                 \label{eq:exptail}
\end{equation}
Indeed the term at degree \(2S\) is bounded by \(4^S\), and all
following term ratios are at most \(1/2\).
To obtain this bound before using exponential addition, the finite
multinomial identity gives
\[
 \sum_{n=0}^m\frac{S^n}{n!}
       \le\left(\sum_{j=0}^m\frac1{j!}\right)^S<3^S
       \qquad(S\text{ a positive integer}).
\]
This argument is entirely finite and removes a dependency on an
already established exponential product law.
Taking \(m=4S+p+4\) supplies more than the requested \(2^{-p}\)
accuracy. The elementary inequality \(\exp1<3<4\) follows directly
from the factorial series. Addition and product identities follow
by finite polynomial multiplication, with the same tail estimate.
For bounded inputs, the useful perturbation bound is
\begin{equation}
 |\exp z-\exp w|\le |z-w|\exp(\max(|z|,|w|)+|z-w|).
 \label{eq:expperturb}
\end{equation}
It follows by factoring the exponential and using the series bound for
\(\exp u-1\).

For positive rational \(t\), define
\[
 \log t=2\sum_{j=0}^{\infty}
                 \frac{y^{2j+1}}{2j+1},\qquad y=\frac{t-1}{t+1}.
\]
On a specified interval \(1/A\le t\le A\), a finite geometric
estimate gives an explicit PR tail. For \(1\le t\le T\), one usable
choice is \(L=T(u+2T+4)\) terms for error at most \(2^{-u}\):
\begin{equation}
 \left|\log t-2\sum_{j<L}\frac{y^{2j+1}}{2j+1}\right|
 \le 2T^2\exp(-(2L+1)/T)<2^{-u}.
 \label{eq:logtail}
\end{equation}
For \(t<1\), use \(\log t=-\log(1/t)\).
Differentiating finite sums and controlling the geometric tails gives
\((\log t)'=1/t\). The finite calculus argument below then proves
\(\exp(\log t)=t\), the logarithm product rule, and the usual
positive-base power laws.

Define
\begin{equation}
       \pi=4\int_0^1\frac{dx}{1+x^2}.                       \label{eq:pidef}
\end{equation}
The integral is a specified rational quadrature limit. Expanding the
finite geometric identity under the integral gives the alternating
arctangent series, with error at most \(1/(2L+1)\) after \(L\)
terms. Thus \(2^{u+3}\) terms give \(\pi\) to error \(2^{-u}\).
Its first eight terms give \(3<\pi\), and the integrand gives
\(\pi<4\).

Write \(a(t)=\int_0^t(1+x^2)^{-1}\,dx\).
Differentiation and a bounded telescoping argument give
\[
 \exp(2\ii a(t))=\frac{1+\ii t}{1-\ii t}.
\]
Both sides equal \(1\) at zero and obey the same explicitly verified
first-order differential identity; multiplying by the inverse of the
rational factor reduces the verification to derivative zero.
At \(t=1\), this yields \(\exp(\ii\pi/2)=\ii\), hence
\(\exp(2\pi\ii)=1\).
There is therefore no circular use of Gaussian normalization in the
definition or period of the trigonometric functions.

We will use
\[
 |\exp(\ii u)|=1,\qquad |\sin u|\le1 \quad(u\in\R).
\]
These follow from the exponential addition and conjugation identities.
On \(0\le u\le2\), alternating Taylor estimates give
\(\sin u\ge u-u^3/6\ge u/3\).
This proves positivity on \([0,\pi/2]\); reflection proves it on
\([0,\pi]\). Consequently \(\sin''=-\sin\le0\), and the finite
calculus version of the chord inequality gives
\begin{equation}
                  \sin u\ge 2u/\pi\quad(0\le u\le\pi/2).
                  \label{eq:sinchord}
\end{equation}
For clarity, the chord inequality needs no extremum theorem: subtract
the affine chord, compare the finite averages of the decreasing
derivative on the two adjacent subintervals, and telescope. Taylor
remainder bounds pass these finite inequalities to every precision.

A positive square root of a \emph{rational} input is computed by
bisection in a known rational interval, using rational comparisons.
The required number of bisections is explicit. General named inputs
are handled by a continuity estimate with a supplied positive lower
bound; no exact comparison with an arbitrary real is made.

\subsection{Quadrature, calculus, and error composition}

\begin{lemma}[Finite calculus]\label{lem:finitecalculus}
Suppose a specified elementary function on a rational interval
\([a,b]\) has Lipschitz constant at most the supplied number \(D\).
Its left rectangle sum on \(q\) equal intervals differs from its
integral name by at most
\[
              \frac{D(b-a)^2}{2q}.
\]
If each sample has error \(\epsilon\), add \((b-a)\epsilon\).
If a specified function \(f\) has second derivative bounded by \(C\),
then for a subdivision \(a=x_0<\cdots<x_q=b\),
\begin{equation}
 \left|f(b)-f(a)-\sum_{j<q}f'(x_j)(x_{j+1}-x_j)\right|
       \le \frac C2\sum_{j<q}(x_{j+1}-x_j)^2.               \label{eq:telescoping}
\end{equation}
\end{lemma}

\begin{proof}
For the quadrature estimate, compare rectangle sums on a common
refinement. On a cell of width \(h\), the total change is at most
\(D h^2/2\); finite summation gives the bound. This makes the
quadrature sums a name with an explicit modulus, so no general
completeness axiom is needed.

For~\eqref{eq:telescoping}, use the Taylor estimate
\(|f(x+h)-f(x)-hf'(x)|\le Ch^2/2\) on each cell and add.
For the elementary functions used here the Taylor estimate comes
from their series, or from the same common-refinement quadrature
calculation applied twice. In particular \(f'=0\) implies
\(|f(b)-f(a)|\le C(b-a)^2/(2q)\) for every specified \(q\).
Choosing \(q\) from the requested precision proves equality.
\end{proof}

Products and integration by parts follow by applying the finite product
identity on each cell; the omitted cross-terms are bounded by a constant
times \(\sum h^2\). Linear changes of variable are reindexings of
finite rectangle sums. The substitution \(t=1/u\), used below on
compact intervals bounded away from zero, follows by comparing the
image partition with a common refinement; its derivative and second
derivative have numerical bounds. The endpoint errors will always
be displayed before an improper integral is used.

These facts apply to the particular explicit functions in the proof.
They are not assumptions about arbitrary differentiable functions.

Finally, if \(a_j,b_j\) are finitely many complex values, then
\begin{equation}
 \left|\prod_{j=1}^m a_j-\prod_{j=1}^m b_j\right|
 \le \sum_{j=1}^m |a_j-b_j|
          \prod_{\ell<j}|a_\ell|\prod_{\ell>j}|b_\ell|.
 \label{eq:producterror}
\end{equation}
This exact telescoping identity controls every product of
approximants. Finite sums use the triangle inequality and precision
increased by the binary length of the number of terms.

\section{A finite proof of the theta transformation}\label{sec:theta}

Fix an integer \(B\ge2\) and initially a rational
\(B^{-1}\le t\le B\). Define
\begin{equation}
 g_t(x)=\exp(-\pi t x^2),\quad
 I(t)=\int_\R g_t(x)\,dx,\quad
 F_t(\xi)=\int_\R g_t(x)\exp(-2\pi\ii\xi x)\,dx.
 \label{eq:Fourierdef}
\end{equation}
The value of \(I(t)\) is initially unknown.

\subsection{Tails and Fourier quadrature}

For \(|x|\ge1\), \(g_t(x)\le\exp(-|x|/B)\); for all \(x\),
\[
 g_t(x)\le2\exp(-|x|/B).
\]
Hence for \(R\ge1\) and every integer \(r\ge0\),
\begin{align}
 \int_{|x|>R}g_t(x)\,dx
       &\le2B\exp(-R/B),                                  \label{eq:gtail}\\
 \int_\R |x|^r g_t(x)\,dx
       &\le4r!B^{r+1}.                                    \label{eq:moments}
\end{align}
The second bound uses \(r\) finite integrations by parts for the
exponential moment. Its tail has an explicit modulus too:
\begin{equation}
 |x|^r\exp(-|x|/B)
       \le(2B)^r r!\exp(-|x|/(2B)).
 \label{eq:momenttail}
\end{equation}
This controls all differentiated improper integrals below.

For \(|\xi|\le K\), a completely explicit quadrature for \(F_t(\xi)\)
at precision \(p\) is obtained by setting
\begin{equation}
 R=B(p+B+4),\quad D=8(BR+K+1),\quad
 q=2^{p+4}(2R)^2D.
 \label{eq:fouriercutoffs}
\end{equation}
The integrand has derivative bounded by \(D\) on \([-R,R]\).
The tail in~\eqref{eq:gtail} is at most \(2^{-p-3}\), the
rectangle error is at most \(2^{-p-5}\), and sample accuracy
\(2^{-(p+R+5)}\) costs at most \(2^{-p-4}\) after integration.
All cutoffs are PR.

\subsection{Fourier transform with unknown normalization}

\begin{lemma}\label{lem:gaussianfourier}
For rational \(\xi\), and hence for the indicated named inputs,
\begin{equation}
                F_t(\xi)=I(t)\exp(-\pi\xi^2/t).
                \label{eq:gaussianfourier}
\end{equation}
\end{lemma}

\begin{proof}
The estimate
\(|\exp(\ii u)-1-\ii u|\le u^2/2\), together with
\eqref{eq:moments}, justifies differentiation with an explicit error:
\begin{equation}
 F_t'(\xi)=-2\pi\ii\int_\R xg_t(x)\exp(-2\pi\ii\xi x)\,dx.
 \label{eq:Fourierderivative}
\end{equation}
The error in the difference quotient is at most \(256B^3|h|\).
On \([-R,R]\), integrate \(g_t'=-2\pi txg_t\) by parts.
The boundary term is at most \(2\exp(-\pi tR^2)\).
Using the moment-tail bound gives
\[
 \int_\R g_t'(x)\exp(-2\pi\ii\xi x)\,dx
       =2\pi\ii\xi F_t(\xi),
\]
so
\[
                    F_t'(\xi)=-(2\pi\xi/t)F_t(\xi).
\]
Set \(H(\xi)=\exp(\pi\xi^2/t)F_t(\xi)\). The product rule gives
\(H'=0\).

To record the finite content, moments give
\[
 |F_t^{(r)}(\xi)|\le4\cdot8^r r!B^{r+1}.
\]
On \(|\xi|\le K\), a bound for \(|H''|\) is
\begin{equation}
 C=2048B^3(K+1)^2\,2^{8B(K+1)^2}.                         \label{eq:Hsecond}
\end{equation}
Indeed put \(a=\pi/t\le4B\); the first two derivatives of
\(\exp(a\xi^2)\) are bounded by
\(8BK\exp(a\xi^2)\) and
\((8B+64B^2K^2)\exp(a\xi^2)\), respectively.
Combine these with \(|F|\le4B\), \(|F'|\le32B^2\),
\(|F''|\le512B^3\).

On a subdivision of the segment from \(0\) to rational \(\xi\)
into \(q\) equal parts, Lemma~\ref{lem:finitecalculus} gives
\[
                 |H(\xi)-H(0)|\le CK^2/(2q).
\]
Choose \(q\ge C(K+1)^2 2^{p+1}\) at precision \(p\), allowing the
separate quadrature errors already given. This proves the identity
at every precision. Since \(H(0)=I(t)\), the claim follows.
\end{proof}

\subsection{Periodization and finite Fourier recovery}

Let
\[
          G_t(x)=\sum_{n\in\Z}g_t(x+n),\qquad
          \theta(t)=G_t(0)=1+2\ps(t).
\]
For \(x\in[0,1]\), omitting \(|n|>N\), \(N\ge2\), costs at most
\begin{equation}
                4B\,2^{-\lfloor(N-1)/B\rfloor}.             \label{eq:periodtail}
\end{equation}
This follows by grouping the exponential tail in blocks of \(B\)
indices. The same block calculation gives
\begin{equation}
 |G_t(x)|\le M:=8B,\qquad
 |G_t(x)-G_t(y)|\le L|x-y|,\quad L:=128B^3.                \label{eq:periodbounds}
\end{equation}
For the derivative bound use
\[
 |g_t'(u)|\le16B^2\exp(-|u|/(2B)),
\]
then sum the geometric blocks. Periodicity gives the bound also
in circular distance.

The Fourier coefficients are
\begin{equation}
 c_k=\int_0^1G_t(x)\exp(-2\pi\ii kx)\,dx
       =F_t(k)=I(t)\exp(-\pi k^2/t).                       \label{eq:coeffs}
\end{equation}
For the middle equality, first take the finite sum
\(-N\le n\le N\), substitute \(u=x+n\) in each integral,
and join the intervals to \([-N,N+1]\). The factor
\(\exp(2\pi\ii kn)=1\) follows from the elementary period proved
earlier. Equations~\eqref{eq:gtail} and \eqref{eq:periodtail}
control the two discarded tails. This proves the exchange without
using a general theorem on infinite sums of integrals.

For an integer \(m\ge1\), the finite Fej\'er kernel is
\begin{align}
 K_m(y)&=\frac1m\left|\sum_{j=0}^{m-1}\exp(2\pi\ii jy)\right|^2\\
 &=\sum_{|k|<m}\left(1-\frac{|k|}{m}\right)
                         \exp(2\pi\ii ky).
 \label{eq:Fejer}
\end{align}
This is a finite identity. Thus \(K_m\ge0\) and
\(\int_{-1/2}^{1/2}K_m=1\). For \(0<|y|\le1/2\), the finite
geometric-sum identity gives
\begin{equation}
 K_m(y)=\frac{\sin^2(\pi my)}{m\sin^2(\pi y)}
                \le\frac1{my^2}.                         \label{eq:Fejerbound}
\end{equation}
Here \(|\sin(\pi y)|\ge |y|\) follows directly from
\(\sin u\ge u/3\) for \(0\le u\le2\) and \(3\le\pi\le4\).

Split the convolution integral at \(\delta=1/h\), \(h\ge2\).
Equations~\eqref{eq:periodbounds} and \eqref{eq:Fejerbound}
give
\begin{align}
 \left|\int_{-1/2}^{1/2}G_t(-y)K_m(y)\,dy-G_t(0)\right|
 &\le L/h+2M\int_{1/h\le|y|\le1/2}\frac{dy}{my^2}\\
 &\le L/h+4Mh/m.                                         \label{eq:Fejererror}
\end{align}
Expanding the finite kernel and using~\eqref{eq:coeffs}, its
convolution value is \(I(t)W_m(1/t)\), where
\[
 W_m(u)=1+2\sum_{k=1}^{m-1}(1-k/m)\exp(-\pi uk^2).
\]
For \(1/B\le u\le B\), the block calculation
\[
 \sum_{k\ge1}k\exp(-\pi uk^2)
       \le B^2\sum_{j\ge0}(j+1)2^{-j}=4B^2
\]
implies
\[
       0\le\theta(u)-W_m(u)\le8B^2/m.
\]
All infinite expressions in these inequalities have the geometric
tail bounds just displayed; the underlying summations are finite.
Since \(I(t)\le4B\), we have proved the decisive estimate
\begin{equation}
 \boxed{|\theta(t)-I(t)\theta(1/t)|
              \le L/h+(4Mh+32B^3)/m.}                    \label{eq:thetafinite}
\end{equation}
At requested precision \(p\), take
\begin{equation}
 h=2^{p+3}(L+1),\qquad
 m=2^{p+3}(4Mh+32B^3+1).                                 \label{eq:Fejercutoffs}
\end{equation}
The bound is at most \(2^{-p-2}\), leaving room for rational
evaluation errors.

\subsection{Gaussian normalization and theta inversion}

At \(t=1\), equation~\eqref{eq:thetafinite} at every precision
gives \(\theta(1)=I(1)\theta(1)\). Its defining positive series
gives \(\theta(1)\ge1\), so
\[
 |I(1)-1|\le |\theta(1)(I(1)-1)|=0.
\]
Thus
\begin{equation}
                  \int_\R\exp(-\pi x^2)\,dx=1.            \label{eq:gaussiannormal}
\end{equation}
The cancellation has the explicit lower bound \(1\).
There is no test for the zero of an arbitrary real.

Approximate the positive square root of rational \(t\) by bisection.
The linear substitution \(y=\sqrt t\,x\), first on finite intervals
and then with~\eqref{eq:gtail}, gives \(I(t)=t^{-1/2}\).
Approximating the scaling factor by rationals uses its positive
lower bound \(1/B\) and the given derivative bounds.
Consequently
\begin{equation}
             \boxed{\theta(t)=t^{-1/2}\theta(1/t).}         \label{eq:theta}
\end{equation}

For later use, this proof is uniform in the parameter.
The block estimate
\[
 \sum_{n\ge1}n^2\exp(-\pi tn^2)\le12B^3
\]
gives a Lipschitz bound \(96B^3\) for \(\theta\) on
\([1/B,B]\). The difference of the two sides
of~\eqref{eq:theta} has Lipschitz bound \(200B^6\).
This extends the rational-input theorem to every specified PR
real name on this interval. Approximants may be clipped to the
larger interval \([1/(2B),2B]\); all clipping comparisons are rational.

\section{Reciprocal gamma from finite products}\label{sec:gamma}

The gamma identities used in the functional equation can be proved
globally without deciding whether a point is a pole. We work with
\(\gamrec_N\) from~\eqref{eq:R}. The classical formulas and their
normalizations are recorded in \cite{DLMF5,DLMF58}; we prove the
needed quantitative versions here.

\subsection{Uniform convergence with an explicit modulus}

For an integer \(B\ge1\), put
\begin{equation}
 M_B=B\,4^{B+B^2},\quad
 C_B=2M_B(B+B^2),\quad
 J_B=2(B+B^2).                                            \label{eq:gammaConstants}
\end{equation}

\begin{lemma}\label{lem:Rconverges}
On \(|z|\le B\), for \(M\ge N\ge J_B\),
\begin{equation}
 |\gamrec_N(z)|\le M_B,\qquad
 |\gamrec_M(z)-\gamrec_N(z)|\le C_B/N.
 \label{eq:Rtail}
\end{equation}
In particular \(|\gamrec(z)-\gamrec_N(z)|\le C_B/N\), and
\(N\ge\max(J_B,2^pC_B)\) is a PR convergence modulus.
\end{lemma}

\begin{proof}
Write \(H_N=\sum_{n=1}^N1/n\) and \(c_N=H_N-\log N\).
Rectangle comparison with the integral of \(1/x\) gives
\[
 0<c_N\le1,\qquad |c_M-c_N|\le1/N.
\]
The elementary inequality
\[
 |(1+w)\exp(-w)|\le\exp(|w|^2/2)
\]
follows by squaring and applying \(1+u\le\exp u\) to
\(u=2\Re w+|w|^2\). Since \(\sum_{n\ge1}n^{-2}\le2\),
\[
 |\gamrec_N(z)|\le B\exp(B+B^2)\le M_B.
\]
For \(n>N\ge2B\), the power series for \(\log(1+z/n)\) gives
\[
 |\log(1+z/n)-z/n|\le B^2/n^2.
\]
Only this local logarithm series is used. Direct multiplication gives
the exact identity
\[
 \gamrec_M(z)=\gamrec_N(z)\exp E,\qquad
 E=z(c_M-c_N)+\sum_{n=N+1}^M(\log(1+z/n)-z/n).
\]
It does not divide by \(\gamrec_N\).
Now \(|E|\le(B+B^2)/N\le1/2\), and
\(|\exp E-1|\le2|E|\); the required bound follows.
\end{proof}

This also gives a direct local power-series description of \(\gamrec\).
Write \(J=J_B\) and \(c=\lim_N(H_N-\log N)\), using the \(1/N\)
modulus just proved. On the disk \(|z|\le B\),
\begin{equation}
 \gamrec(z)=z\prod_{n=1}^{J}(1+z/n)\,
 \exp\left(z(c-H_J)+
        \sum_{n>J}\bigl(\log(1+z/n)-z/n\bigr)\right).
 \label{eq:Rlocalfactor}
\end{equation}
The differentiated tails are bounded geometric or reciprocal-square
series. Thus \(\gamrec\) is entire with effective local approximants.
The exponential has the explicit inverse obtained by negating its
exponent. This finite-polynomial-times-exponential factorization
also shows that the denominators used later represent nonzero
analytic germs; their possible zeros come from the listed linear factors.

Rearranging the finite products gives, for every \(z\),
\begin{equation}
 z\gamrec_N(z+1)=\left(1+\frac{z+1}{N}\right)\gamrec_N(z).
\end{equation}
The controlled limit yields
\begin{equation}
 \gamrec(z)=z\gamrec(z+1),\qquad \gamrec(0)=0,\quad\gamrec(1)=1.
 \label{eq:Rrecurrence}
\end{equation}
These are entire identities, including their zero arguments.

\subsection{Identification with Euler's gamma integral}

For \(\Re z=\sigma\ge1/A>0\), \(|z|\le B\), define
\[
                  \Gamma(z)=\int_0^\infty \exp(-u)u^{z-1}\,du.
\]
The lower endpoint error at \(\epsilon\le1\) is at most
\(A\epsilon^{1/A}\). The tail above \(T\ge1\) is at most
\begin{equation}
                  2^{B+1}B!\exp(-T/2),                   \label{eq:gammaintegraltail}
\end{equation}
using \(u^B\le2^BB!\exp(u/2)\). Quadrature on
\([\epsilon,T]\) has a derivative bound obtained by differentiating
this explicit integrand. Thus the integral has a PR name once its
positive real-part margin is supplied.

Let
\[
 \Gamma_N(z)=\int_0^N u^{z-1}(1-u/N)^N\,du.
\]
Substitute \(u=Nv\), and integrate the beta polynomial by parts \(N\)
times. The endpoint \(0\) is controlled by the same real-part margin.
The result is the finite identity
\begin{equation}
 \Gamma_N(z)=
 \frac{N^zN!}{z(z+1)\cdots(z+N)},\qquad
                  \Gamma_N(z)\gamrec_N(z)=1.              \label{eq:betaEuler}
\end{equation}
All denominators have positive real part.

For \(0\le u\le N/2\),
\begin{equation}
 0\le \exp(-u)-(1-u/N)^N
                  \le (u^2/N)\exp(-u).                   \label{eq:betakernel}
\end{equation}
Indeed the series for \(-N\log(1-u/N)-u\) is at most
\(u^2/[2N(1-u/N)]\le u^2/N\), and
\(1-\exp(-v)\le v\). On the rest of the interval, both nonnegative
kernels are bounded by \(\exp(-u)\).
Integrating these estimates gives, for \(N\ge2\),
\begin{equation}
 |\Gamma_N(z)-\Gamma(z)|\le D_B/N,\qquad
 D_B=1+(B+1)!+2^{B+3}B!.                                 \label{eq:gammabetaerror}
\end{equation}
For the first part, split at \(u=1\) and use
\(u^{\sigma+1}\le1+u^{B+1}\).
For the tail use~\eqref{eq:gammaintegraltail} at \(N/2\)
and \(\exp(-N/4)\le4/N\).

Moreover \(|\Gamma_N(z)|\le A+B!\).
Combining~\eqref{eq:Rtail}, \eqref{eq:betaEuler},
and~\eqref{eq:gammabetaerror}, with the finite product estimate,
proves
\begin{equation}
                 \Gamma(z)\gamrec(z)=1\qquad(\Re z>0).
                 \label{eq:gammaReciprocal}
\end{equation}
Thus our reciprocal-gamma construction agrees with Euler's integral
on its defining half-plane. No analytic continuation principle has
been used in this identification.

\subsection{Euler's sine product by finite trigonometric algebra}

We need the global identity
\begin{equation}
 \frac{\sin(\pi z)}{\pi}
       =z\prod_{n=1}^\infty(1-z^2/n^2).                  \label{eq:sineproduct}
\end{equation}
Its classical normalization is \cite{DLMF422}. The following finite
argument also provides a PR modulus.

Let \(N=2m+1\) be odd and \(w=\pi z/N\).
The addition formulas give
\[
 \sin(Nw)=\sin w\,F_N(\sin^2w),\quad
 F_N(Y)=\sum_{j=0}^m(-1)^j\binom N{2j+1}Y^j(1-Y)^{m-j}.
\]
The \(m\) values \(\sin^2(\pi k/N)\), \(1\le k\le m\), are
roots of \(F_N\). They are separated: for \(j<k\),
\[
 \sin^2(\pi k/N)-\sin^2(\pi j/N)
  =\sin(\pi(k+j)/N)\sin(\pi(k-j)/N)\ge4/N^2.
\]
Use~\eqref{eq:sinchord} and reflection about \(\pi/2\).
Their denominators are therefore explicitly nonzero.
Finite polynomial division, together with \(F_N(0)=N\), proves
\begin{equation}
 \sin(\pi z)=N\sin(\pi z/N)
   \prod_{k=1}^m
   \left(1-\frac{\sin^2(\pi z/N)}{\sin^2(\pi k/N)}\right).
 \label{eq:finitesine}
\end{equation}
No root of an unspecified polynomial is selected.

Here are sufficient quantitative bounds for passing to the product.
Suppose \(|z|\le B\), \(B\ge1\), and \(N\ge8B^2\). Set
\[
 a_k=\frac{\sin^2(\pi z/N)}{\sin^2(\pi k/N)},\qquad
 b_k=\frac{z^2}{k^2}.
\]
For \(0<v\le\pi/2\),
\[
 0\le\frac1{\sin^2v}-\frac1{v^2}\le2.
\]
In fact \(v-\sin v\le v^3/6\), \(\sin v\le v\), and
\(\sin v\ge v/2\) bound the left side by \(4/3\).
Taylor estimates at \(|\pi z/N|\le1/2\) now give
\begin{align}
 \sum_{k=1}^m|a_k-b_k|
     &\le128B^4/N^2+16B^2/N\le32B^2/N,\\
 \sum_{k=1}^m|a_k|&\le32B^2,\qquad
 \sum_{k=1}^m|b_k|\le2B^2,\\
 |N\sin(\pi z/N)/\pi-z|&\le6B^3/N^2.
 \label{eq:sinebounds}
\end{align}
Formula~\eqref{eq:producterror} bounds the product discrepancy by
\(32B^2\,4^{34B^2}/N\).
For \(P_m(z)=\prod_{k=1}^m(1-z^2/k^2)\), the reciprocal-square
tail gives
\[
 |P_\ell(z)-P_m(z)|
          \le2\exp(2B^2)B^2/m\quad(\ell\ge m\ge B^2).
\]
The proof is the same finite telescoping formula: all mixed partial
products are at most \(\exp(2B^2)\), and
\(\sum_{k>m}k^{-2}\le1/m\).

Since \(m\ge N/3\), the finite prefactor error, product comparison,
and product tail in~\eqref{eq:finitesine} have total bound
\begin{equation}
                  \frac{C_{\sin}(B)}{N},\qquad
 C_{\sin}(B)=44B^3\,4^{34B^2}.                            \label{eq:sinemodulus}
\end{equation}
Their three contributions in these common units are at most
\(6+32+6\). Choose the next odd integer above
\(\max(8B^2,2^pC_{\sin}(B))\).
This proves~\eqref{eq:sineproduct} at every prescribed precision,
including \(z=0\), with no division by \(z\).

\subsection{Reflection, duplication, and the required factor identity}

Direct rearrangement of finite products yields
\begin{align}
 \gamrec_N(z)\gamrec_N(1-z)
 &=z\left(1+\frac{1-z}{N}\right)
               \prod_{n=1}^N(1-z^2/n^2),                 \label{eq:finiteReflection}\\
 \gamrec_N(z)\gamrec_N(z+1/2)
 &=2^{2z-1}\left(1+\frac{z}{N+1/2}\right)
                    \gamrec_N(1/2)\gamrec_{2N}(2z).
 \label{eq:finiteDuplication}
\end{align}
These are finite polynomial identities after combining the
positive-base powers. To check duplication explicitly, split the
product for \(\gamrec_{2N}(2z)\) into its even and odd indices
and use
\[
 a_N=\prod_{k=1}^N\frac{k-1/2}{k},\qquad
 \gamrec_N(1/2)=\frac{N+1/2}{\sqrt N}\,a_N.
\]
All remaining denominators are positive integers or half-integers.

The established moduli give reflection:
\begin{equation}
       \gamrec(z)\gamrec(1-z)=\sin(\pi z)/\pi.              \label{eq:Rreflection}
\end{equation}
At \(z=1/2\), the finite products are positive, so
\(\gamrec(1/2)\ge0\) and
\(\gamrec(1/2)^2=1/\pi\).
The positive square root, with \(3<\pi<4\), gives
\(\gamrec(1/2)=1/\sqrt\pi\). Consequently duplication is
\begin{equation}
 \gamrec(z)\gamrec(z+1/2)
             =\frac{2^{2z-1}}{\sqrt\pi}\gamrec(2z).
 \label{eq:Rduplication}
\end{equation}

Apply~\eqref{eq:Rduplication} at \(z=(1-s)/2\), rearranging only
the everywhere-apart exponential factor:
\[
 \gamrec(1-s)=2^s\sqrt\pi\,
                  \gamrec((1-s)/2)\gamrec(1-s/2).
\]
Multiply by \(\gamrec(s/2)\) and use reflection.
We obtain the entire identity
\begin{equation}
 \boxed{\gamrec(1-s)\gamrec(s/2)
 =\frac{2^s}{\sqrt\pi}\sin(\pi s/2)\gamrec((1-s)/2).}
 \label{eq:gammafactor}
\end{equation}
No potentially zero gamma factor has been cancelled.

\subsection{Explicit apartness when a point value is wanted}

The entire identities above need no pole exclusion. To evaluate
\(\Gamma(z)\) as a number, suppose \(|z|\le B\), set \(J=J_B\), and
supply a rational \(0<\delta\le1\) with
\[
                   |z+j|\ge\delta\quad(0\le j\le J).
\]
The finite initial product and the tail-exponential bound give
\begin{equation}
                |\gamrec(z)|\ge
                  \frac{\delta^{J+1}}{2J!J^B}.             \label{eq:Rlower}
\end{equation}
Indeed the initial factors have this lower bound without the \(2\);
the remaining exponent has modulus at most \(1/2\).
This supplies a PR reciprocal modulus. It also explains, with
explicit witnesses, the possible poles at nonpositive integers.
There is no need to decide whether an arbitrary named input equals
one of them.

\section{The regularized Mellin construction}\label{sec:mellin}

We now prove that the symmetric expression \(X\) really continues the
completed Dirichlet series. This identification is the substantive
content behind the simple symmetry in~\eqref{eq:X}.

\subsection{Uniform bounds and a fully specified finite approximation}

For \(t\ge1\), since \(n^2\ge1+3(n-1)\),
\begin{equation}
 \ps(t)\le\frac{\exp(-3t)}{1-\exp(-9t)}
                    \le2\exp(-3t).                      \label{eq:psitail}
\end{equation}
For an integer \(B\ge2\) and \(|\Re s|,|\Im s|\le B\),
\(|t^{s/2-1}|\le t^B\). The factorial-series inequality
\(t^B\le B!\exp t\) gives
\begin{equation}
 \normone{\int_T^\infty\ps(t)t^{s/2-1}\,dt}
             \le2B!\exp(-2T)\le2B!2^{-2T}.                \label{eq:Jtail}
\end{equation}
In particular \(\normone{J(s)}\le2B!\).

For a nonnegative integer \(n\), let
\[
 \lceilog(0)=0,\qquad
 \lceilog(n)=\min\{j\le n:n\le2^j\}\quad(n\ge1).
\]
This is a bounded PR search, the ceiling binary logarithm.
At precision \(p\), set
\begin{align}
 T&=p+\lceilog(B!)+4,&
 N&=p+(B+1)\lceilog(T)+4,\nonumber\\
 D&=2(B+2)T^B,&
 q&=p+\lceilog(TD)+4,\nonumber\\
 r&=p+\lceilog(T)+4,&
 M&=(T-1)2^q.                                            \label{eq:Jcutoffs}
\end{align}
Let \(t_j=1+j2^{-q}\), \(0\le j<M\), and
\begin{equation}
 g_N(t,s)=\sum_{n=1}^N
        \exp\bigl(-\pi n^2t+(s/2-1)\log t\bigr).
 \label{eq:gn}
\end{equation}
Approximate this finite sum by a rational complex value
\(\widehat g_N(t,s)\) with \(\ell^1\)-error at most \(2^{-r}\).
It suffices to approximate each of its \(N\) terms to precision
\(r+\lceilog(N)+1\), leaving half the sample budget for rounding.
Define
\begin{equation}
            \Jh(B,p,s)=2^{-q}\sum_{j=0}^{M-1}
                                    \widehat g_N(t_j,s).  \label{eq:Jhat}
\end{equation}
This is an explicit finite rational expression.

\begin{proposition}[Mellin approximation contract]\label{prop:Jcontract}
For every valid rational input in the stated box,
\[
                  \normone{\Jh(B,p,s)-J(s)}<2^{-p}.
\]
More precisely, a sufficient error budget is
\begin{equation}
 2B!2^{-2T}
 +4T^{B+1}2^{-3(N+1)^2}
 +TD2^{-q}
 +T2^{-r}.                                               \label{eq:Jbudget}
\end{equation}
\end{proposition}

\begin{proof}
The first term is~\eqref{eq:Jtail}.
For \(n=N+1+j\), \(n^2\ge(N+1)^2+j\).
Thus the omitted Gaussian series on \(t\ge1\) is bounded by
\(2\cdot2^{-3(N+1)^2}\). The complex weight has
\(\ell^1\)-norm at most \(2T^B\); integration over length at most
\(T\) gives the second term.

For \(\ps_N(t)=\sum_{n=1}^N\exp(-\pi n^2t)\), \(t\ge1\),
we have \(\ps_N\le1\) and \(|\ps_N'|\le1\).
For the latter, use \(n^2\le4^{n-1}\), \(n^2\ge n\), and
\[
       4n^2\,2^{-3n^2}\le2^{-n}.
\]
Then \(|s/2-1|\le B+1\) and differentiation of the power give
\(\normone{g_N'(t,s)}\le2(B+2)T^B=D\).
Lemma~\ref{lem:finitecalculus} bounds the mesh error by
\(TD2^{-q}\). Summing the sample errors contributes \(T2^{-r}\).

For completeness, divide~\eqref{eq:Jbudget} by \(2^{-p}\).
Its four components are at most
\[
                    \frac1{128},\quad\frac1{32},
                    \quad\frac1{16},\quad\frac1{16}.
\]
For the second bound put \(d=(B+1)\lceilog(T)\);
then \(N=p+d+4\), and its ratio is at most
\[
             2^{\,N-2-3(N+1)^2}\le1/32.
\]
The other bounds follow immediately from~\eqref{eq:Jcutoffs}
and \(B!\le2^{\lceilog(B!)}\).
Their sum is \(21/128<1\).
\end{proof}

\paragraph{A completely bounded sample evaluator.}
The sample accuracy used above does not rely on waiting for convergence.
Suppose \(1\le t\le T\), \(n\le N\), and the required term accuracy is
\(2^{-u}\). Put
\begin{equation}
 U=(4N^2+B+2)T,\qquad C=N^2T+B+1,\qquad
 v=u+2U+C+4.
 \label{eq:sampleprecision}
\end{equation}
Compute \(\pi\) and \(\log t\) to error \(2^{-v}\) by the
bounded algorithms in Section~\ref{sec:toolkit}. Form the
exponent rationally and compute its exponential to precision
\(u+1\) by~\eqref{eq:exptail}.
The true exponent has \(\ell^1\)-norm at most \(U\); its input error
is at most \(C2^{-v}\le2^{-u-2U-4}\).
The exponential perturbation costs at most \(2^{-u-1}\), using
the safe factor \(2^{2U+3}\); the final Taylor error costs at most
another \(2^{-u-1}\). This proves the advertised sample contract.

The bounds are intentionally generous. For example \(B=4,p=20\)
gives \(T=29,N=49,q=52\) and
\(M=126100789566373888\).
This is a primitive recursive existence and proof-theoretic
construction, not a claim of a practical zeta evaluator. The numerical
diagnostics use a separate fast method and are labelled accordingly.

\subsection{Local power series and named parameters}

The integral defining \(J\) is entire in a directly effective sense.
For \(h\) in a bounded disk, expand
\[
 t^{(s+h)/2-1}
       =t^{s/2-1}\sum_{j=0}^\infty
                         \frac{(h\log t/2)^j}{j!}.
\]
Here is a precise finite justification. Choose an integer \(H_0\)
with \(|h|\le H_0\) on the whole disk, and use the enlarged
box \(B'=B+H_0+2\). Choose \(T\) so that twice the
tail in~\eqref{eq:Jtail} for \(B'\) is within the allotted error.
On \([1,T]\), \(|h\log t/2|\) is explicitly bounded.
The exponential remainder~\eqref{eq:exptail} makes the integrated
Taylor error arbitrarily small with a PR degree. On the omitted
tail, the exponential majorant of the sum is at most
\(t^{|h|/2}\); the same enlarged-box bound applies both to
the original integral and to the finite Taylor polynomial whose
coefficient integrals extend to infinity.
Finite sums and quadrature may be interchanged exactly before
these two errors are applied. This proves a local convergent
power series with coefficients
\[
 \frac1{2^j j!}\int_1^\infty
          \ps(t)t^{s/2-1}(\log t)^j\,dt.
\]
All coefficients have effective tails because \(\log t\le t\).
For example the \(j\)-th differentiated tail is bounded by
\(2(B+j)!2^{-2T}\) on an appropriate original box.

The same calculation supplies a numerical local continuity bound:
along a segment in a box of radius \(B\),
\[
 \normone{J(s)-J(w)}\le (B+1)!\,|s-w|
\]
is sufficient. Indeed the parameter derivative has integrand
\(\frac12\ps(t)t^{s/2-1}\log t\), and the preceding factorial
majorant applies. Enlarging the box absorbs rounding of the input.
This permits evaluation on fixed PR names by reading sufficiently
many input digits. Uniform derivative bounds for \(\gamrec\)
come from its polynomial-times-exponential description in
Section~\ref{sec:gamma}; alternatively each finite evaluator has
a directly computable sensitivity bound, and its uniform tail
handles the passage to the name.

\subsection{The Mellin identity on the convergence half-plane}

Assume now that \(\sigma=\Re s\ge1+\delta\), where a positive
rational margin \(\delta\) is supplied. On \(\sigma>1\), define
the Dirichlet-series name
\[
                         \zeta_D(s)=\sum_{n=1}^\infty n^{-s}.
\]
Its tail after \(N\) is at most
\begin{equation}
               \sum_{n>N}n^{-\sigma}
                    \le N^{1-\sigma}/(\sigma-1)
                    \le\delta^{-1}N^{-\delta}.            \label{eq:Dirichlettail}
\end{equation}
This follows by comparing the decreasing positive terms with
finite integrals of \(x^{-\sigma}\), then using their explicit
endpoint values. For \(\delta\ge1/a\), \(a\ge1\), an exponential
cutoff \(N=2^{a(p+a+c)}\), with \(c\) a binary-size allowance for
the displayed constant, is sufficient at precision \(p\).
Thus the tail is PR even for a very small supplied margin.

For each fixed positive integer \(n\), the linear substitution
\(u=\pi n^2t\) gives
\begin{equation}
 \int_0^\infty \exp(-\pi n^2t)t^{s/2-1}\,dt
                  =\pi^{-s/2}n^{-s}\Gamma(s/2).
 \label{eq:singleMellin}
\end{equation}
The gamma endpoint bounds make the substitution effective.
For a finite number of \(n\)'s we can sum this identity.
We next justify its infinite version quantitatively.

The Gaussian normalization and monotonicity give
\[
 \ps(t)\le\int_0^\infty \exp(-\pi tx^2)\,dx
                           =\frac1{2\sqrt t}.
\]
Consequently the Mellin lower endpoint below \(\epsilon\) has
error at most
\begin{equation}
 \int_0^\epsilon \ps(t)t^{\sigma/2-1}\,dt
                \le\frac{\epsilon^{(\sigma-1)/2}}{\sigma-1}.
 \label{eq:Mellinlower}
\end{equation}
The upper endpoint is bounded by~\eqref{eq:Jtail}.
If also \(\sigma\le B\), the positive integrals for the omitted
\(n\)'s have total bound
\begin{equation}
 (2+B!)\,\delta^{-1}N^{-\delta}.                           \label{eq:Mellinnsumtail}
\end{equation}
Indeed \(\Gamma(\sigma/2)\le2+B!\), by splitting its integral
at \(1\), and \(\pi^{-\sigma/2}\le1\). Apply the finite
identity~\eqref{eq:singleMellin} to any finite block \(n>N\)
and then~\eqref{eq:Dirichlettail}.

To make the exchange fully explicit, choose an integer \(a\ge1\)
with \(\delta\ge1/a\), and put
\[
 c=\lceilog(2+B!)+4,\qquad
 N=2^{a(p+a+c)},\qquad
 \epsilon=2^{-2a(p+a+4)}.
\]
The lower endpoint and omitted-\(n\) integral errors are then each
at most \(2^{-p-4}\). Choose \(T\) from~\eqref{eq:Jcutoffs},
with extra accuracy if needed for the remaining budget.
Uniform Gaussian convergence on \([\epsilon,T]\) follows from
the theta estimates with a parameter envelope at least
\(\max(T,1/\epsilon)\).
On this compact interval exchange only a finite sum with a finite
rectangle quadrature. To identify the residual \(\ps\) sum,
first use a finite block \(N<n\le M\) and then the uniform tail
bound as \(M\) is increased to its explicit cutoff.
The three discarded errors are smaller than the prescribed
accuracy. This proves
\begin{equation}
 \boxed{\pi^{-s/2}\Gamma(s/2)\zeta_D(s)
        =\int_0^\infty\ps(t)t^{s/2-1}\,dt
          \qquad(\Re s>1).}                              \label{eq:Mellin}
\end{equation}
No general dominated-convergence or Fubini theorem is assumed.

\subsection{Splitting the integral and constructing the continuation}

The theta equation~\eqref{eq:theta} is equivalent to
\begin{equation}
 \ps(t)=t^{-1/2}\ps(1/t)+\tfrac12(t^{-1/2}-1).
 \label{eq:psitransform}
\end{equation}
On \(\Re s>1\), split~\eqref{eq:Mellin} at \(1\) and substitute
this formula below \(1\). The elementary terms integrate to
\begin{equation}
 \frac12\int_0^1
        \left(t^{(s-1)/2-1}-t^{s/2-1}\right)dt
              =\frac1{s-1}-\frac1s=\frac1{s(s-1)}.
 \label{eq:elementaryMellin}
\end{equation}
At a lower cutoff \(\epsilon\), these equalities follow by
differentiating \(t^a/a\). Their endpoint terms tend to zero with
the positive real-part margin already supplied.
In the remaining integral substitute \(u=1/t\):
\begin{equation}
 \int_0^1 t^{s/2-3/2}\ps(1/t)\,dt
       =\int_1^\infty u^{-s/2-1/2}\ps(u)\,du=J(1-s).
 \label{eq:inverseMellin}
\end{equation}
Use the substitution first on \([\epsilon,1]\).
The omitted transformed interval is a tail of \(J(1-s)\),
so~\eqref{eq:Jtail} gives its modulus.
Although the theta proof uses a parameter box whose bound grows
as \(1/\epsilon\), that bound is itself PR in the chosen precision.

Combining the finite identities gives
\begin{equation}
 \pi^{-s/2}\Gamma(s/2)\zeta_D(s)
        =\frac1{s(s-1)}+J(s)+J(1-s)
                    \qquad(\Re s>1).                     \label{eq:regularizedMellin}
\end{equation}
Multiplication by \(s(s-1)/2\) yields exactly
\begin{equation}
 X(s)=\tfrac12s(s-1)\pi^{-s/2}\Gamma(s/2)\zeta_D(s).
 \label{eq:xiIdentification}
\end{equation}
The right side has now been identified with the entire, directly
specified function \(X\) in~\eqref{eq:X}.

Since \(\gamrec(s/2)=(s/2)\gamrec(1+s/2)\) and
\(\Gamma(s/2)\gamrec(s/2)=1\), multiplication gives
\begin{equation}
                  Z(s)=(s-1)\zeta_D(s)\qquad(\Re s>1).
 \label{eq:ZIdentification}
\end{equation}
The expression~\eqref{eq:Z} is entire everywhere, so
\(\zeta=(Z,s-1)\) is the required explicitly constructed
meromorphic continuation of the Dirichlet series.

\section{The functional equation and all exceptional points}\label{sec:FE}

The elementary identities
\[
 (1-s)((1-s)-1)=s(s-1),\qquad 1-(1-s)=s
\]
show immediately from~\eqref{eq:X} that \(X(1-s)=X(s)\).
It is crucial that Section~\ref{sec:mellin} has already identified
this function with the completed Dirichlet series.

For the sine--gamma form, start with
\[
 s\gamrec(1-s)Z(s)
       =s\gamrec(1-s)\pi^{s/2}\gamrec(1+s/2)X(s).
\]
Use recurrence and then~\eqref{eq:gammafactor}:
\begin{align}
 s\gamrec(1-s)Z(s)
 &=2\pi^{s/2}\gamrec(1-s)\gamrec(s/2)X(s)\\
 &=2^{s+1}\pi^{(s-1)/2}
                  \sin(\pi s/2)\gamrec((1-s)/2)X(s).
 \label{eq:FEleft}
\end{align}
On the other hand, recurrence at \((1-s)/2\) and symmetry give
\begin{align}
 &(1-s)2^s\pi^{s-1}\sin(\pi s/2)Z(1-s)\\
 &\quad=(1-s)2^s\pi^{s-1+(1-s)/2}\sin(\pi s/2)
                    \gamrec(1+(1-s)/2)X(1-s)\\
 &\quad=2^{s+1}\pi^{(s-1)/2}
                    \sin(\pi s/2)\gamrec((1-s)/2)X(s).
 \label{eq:FEright}
\end{align}
The last expressions are identical. This proves the global entire
identity~\eqref{eq:entireFE} by multiplication and addition alone.

\subsection{Meromorphic interpretation}

For this report, a meromorphic expression is represented locally by a
pair of explicit analytic numerator and denominator names. An identity
is proved by equality of their cross-products. Point evaluation of a
quotient is performed when a positive rational lower bound for its
denominator is supplied. This is sufficient for all formulas here;
we do not need a general quotient construction for arbitrary meromorphic
functions.

The left side of~\eqref{eq:standardFE} has pair
\((Z(s),s-1)\).
The right side has pair
\[
 \left(2^s\pi^{s-1}\sin(\pi s/2)Z(1-s),
                              -s\gamrec(1-s)\right).
\]
Their cross-product equality is precisely~\eqref{eq:entireFE}.
Therefore the displayed sine--gamma functional equation is proved
as an identity of meromorphic expressions on the whole plane.
On a region with the appropriate apartness witnesses, it gives
the ordinary equality of numerical values by division.

Similarly the completed function is represented by
\begin{equation}
 \Lambda(s)=
  \left(1+s(s-1)(J(s)+J(1-s)),\;s(s-1)\right).
 \label{eq:Lambda}
\end{equation}
The numerator is \(2X(s)\) and is invariant under \(s\mapsto1-s\);
the denominator has the same invariance.
On \(\Re s>1\), it is
\(\pi^{-s/2}\Gamma(s/2)\zeta_D(s)\).

\subsection{Values, residue, and removable factors}

From~\eqref{eq:X},
\[
 X(0)=X(1)=\frac12.
\]
The reciprocal-gamma formulas give
\[
 \gamrec(1)=1,\qquad
 \gamrec(1/2)=1/\sqrt\pi,\qquad
 \gamrec(3/2)=2/\sqrt\pi.
\]
It follows that
\[
 Z(0)=\frac12,\qquad Z(1)=1.
\]
Hence
\begin{equation}
 \zeta(0)=-\frac12,\qquad
                  \mathop{\rm Res}_{s=1}\zeta(s)=1.       \label{eq:specialvalues}
\end{equation}
In local power-series language, \(Z(s)-1\) has a factor \(s-1\);
so \(\zeta(s)=1/(s-1)\) plus an analytic function near \(1\).
This proves that the pole is simple and genuine.

For every positive integer \(m\), recurrence gives
\(\gamrec(1-m)=0\). Thus \(Z(-2m)=0\) and
\(\zeta(-2m)=0\).
The entire identity~\eqref{eq:entireFE} remains valid at
\(s=0,1,2,\ldots\); it never forms an undefined product.
For example, at a positive even integer the displayed sine and
gamma factors in~\eqref{eq:standardFE} must be interpreted together
meromorphically. Their separate totalized computer values are not
the assertion being proved.

\section{Eliminating the analysis inside iPRA}\label{sec:arithmetic}

We now specify the internal statement and justify the logical strength
of the preceding human proof. All infinite notation has served to
organize finite estimates with explicit moduli.

\subsection{One total rational residual checker}

Use natural codes for signed integers, positive-denominator rationals,
and pairs of rationals. Arithmetic can use canonical fractions, or
cross-product equality with a canonical final form. Integer
division, gcd, bit length, finite sums, factorial, bounded
products, and bounded search are PR. Malformed codes are assigned
fixed defaults.

For valid \(s=x+\ii y\), \(|x|,|y|\le B\), \(B\ge2\), and
precision \(k\), put
\begin{equation}
       A=B+3,\qquad H=2^{64A^2},\qquad
       w=k+512A^2+20.                                    \label{eq:workingprecision}
\end{equation}
Expand the two sides of the entire equation as
\begin{align}
 L(s)&=s\gamrec(1-s)\pi^{s/2}\gamrec(1+s/2)X(s),\nonumber\\
 U(s)&=(1-s)2^s\pi^{s-1}\sin(\pi s/2)
       \pi^{(1-s)/2}\gamrec(1+(1-s)/2)X(1-s).
 \label{eq:expandedFE}
\end{align}
Every factor in these expressions has \(\ell^1\)-norm at most \(H\).
For gamma use~\eqref{eq:gammaConstants} with disk bound \(2A\).
For \(J\) use \(\normone{J}\le2A!\), giving
\(\normone X\le1/2+8A^2A!\le2^{8A^2}\).
The elementary powers and sine are bounded by their exponential
majorants. These estimates are all well below \(H\).

Evaluate every analytic factor to \(\ell^1\)-error at most \(2^{-w}\),
and keep rational multiplications and the factors \(s,1-s\) exact.
The following prescriptions make this a fixed total PR algorithm.

\begin{enumerate}
\item For \(J\), use~\eqref{eq:Jcutoffs}--\eqref{eq:Jhat}.
For \(X\), use the common envelope \(A\), put \(p_J=w+4A+8\), and form
\begin{equation}
 \widehat X(A,w,s)=\frac12+\frac{s(s-1)}2
              \bigl(\Jh(A,p_J,s)+\Jh(A,p_J,1-s)\bigr).
 \label{eq:Xhat}
\end{equation}
Its error is at most \(4A^2 2^{-p_J}<2^{-w}\).
Using the same parameters on both arguments makes
\(\widehat X(A,w,1-s)=\widehat X(A,w,s)\) an exact rational identity.

\item For \(\gamrec(z)\), use \(D=2A\), which bounds the
\(\ell^1\)-norm as well as the modulus of each supplied gamma
argument, and
\[
 N=\max(J_D,\,2^{w+4}C_D).
\]
The tail is at most \(2^{-w-4}\) in complex modulus.
Compute the finite rational product exactly and approximate
\(\exp(-z\log N)\).
One fully bounded precision schedule is
\[
 Q=D(N+1)^D,\quad S=D(N+1)+1,\quad
 v=w+2S+Q+D+14.
\]
The polynomial factor has modulus at most \(Q\):
\(\prod_{n=1}^N(1+D/n)=\binom{N+D}{D}\le(N+1)^D\).
Compute \(\log N\) and the resulting exponential to precision \(v\).
The exponent has modulus at most \(S\), so
\eqref{eq:expperturb} and the factor \(Q\) leave ample
\(\ell^1\) budget for \(2^{-w}\) accuracy.

\item For the remaining powers and sine, fix \(v=w+16A+30\).
Compute \(\pi\), \(\log2\), and
\(\log(\widehat\pi)\) to precision \(v\), first clipping
\(\widehat\pi\) rationally to \([3,4]\). Use the exact finite
series prescribed in Section~\ref{sec:toolkit}.
The error of \(\log(\widehat\pi)\) as an approximation to
\(\log\pi\) is at most \(2\cdot2^{-v}\).
Form each exponent rationally, and evaluate its exponential
to precision \(w+4\). All these exponents have
\(\ell^1\)-norm at most \(4A+1\), their perturbations are at most
\(4A2^{-v}\), and~\eqref{eq:expperturb} proves error below
\(2^{-w}\).
For the sine use
\(\sin z=(\exp(\ii z)-\exp(-\ii z))/(2\ii)\), with the same
exponential accuracy. Use exact finite polynomial sums throughout;
no optional final rounding is part of this fixed checker.
\end{enumerate}

Call the resulting rational products \(\widehat L,\widehat U\).
The telescoping product estimate~\eqref{eq:producterror}, allowing
up to eight factors on either side, bounds their aggregate
evaluation error by
\begin{equation}
                2^{11}H^7\,2^{-w}<2^{-k-8}.               \label{eq:finalbudget}
\end{equation}
All intermediate finite identities in the human proof are used with
an additional finite-precision allowance; their explicit moduli can
make that allowance at most \(2^{-k-8}\) as well.
Equivalently, the analytic shorthand proves \(L=U\), while the finite
estimates just given remove that shorthand at the requested precision.
Define the nontrivial branch of the checker by
\begin{equation}
 \mathsf{Check}(c,B,k)=1
       \quad\Longleftrightarrow\quad
              \normone{\widehat L-\widehat U}\le2^{-k}.
 \label{eq:Checkdefinition}
\end{equation}
Return \(1\) on malformed codes, \(B<2\), or out-of-box inputs.
This specifies one total PR function of three natural arguments.

\subsection{Why the proof needs only quantifier-free induction}

The following list identifies the actual induction invariants.
Each is a rational equation or inequality with all bounds and
precisions supplied as natural parameters.

\begin{longtable}{p{0.22\textwidth}p{0.68\textwidth}}
\toprule
Finite operation & Induction invariant and error accounting\\
\midrule
\endhead
Geometric tails &
The exact finite geometric-sum identity and a bound on the remaining
geometric term. Polynomial weights such as \(j\), \(j^2\), or a
fixed factorial moment are accumulated by a PR fold.\\
Taylor evaluation &
A recurrence for the current rational term and partial sum, together
with the explicit factorial or ratio majorant. The output precision
is fixed throughout this induction.\\
Quadrature &
After \(j\) cells, accumulated error is at most the sum of the first
\(j\) cell budgets. Formula~\eqref{eq:telescoping} is the corresponding
telescoping invariant.\\
Products &
The finite identity~\eqref{eq:producterror}, with a numerical bound
for each mixed product. No induction is applied to an infinite-product
convergence assertion.\\
Fourier algebra &
Finite sums over \(|k|<m\), with their exact geometric identities;
omitted series terms and quadrature errors are bounded by the stated
numerical cutoffs.\\
Gamma and sine identities &
Finite polynomial product rearrangements and polynomial division by
the explicitly separated trigonometric roots; each denominator has
a supplied rational apartness bound.\\
Mellin splitting &
Finite sums and finite rectangle substitutions on
\([\epsilon,T]\); the three endpoint and series tails are
separate rational majorants.\\
\bottomrule
\end{longtable}

A bounded universal assertion, for example that every cell error
among the first \(j\) cells satisfies its bound, can be encoded as
a finite Boolean conjunction produced by a PR function and equated
to \(1\). It is therefore quantifier free. Its elementary relation
to the component tests is itself proved by quantifier-free induction.
Rational inequalities are reduced to natural cross-product
inequalities, with the denominators positive by construction.

For every application of a displayed analytic identity, choose the
requested accuracy symbolically, insert its explicit truncation and
mesh functions, and use the corresponding finite error inequality.
Finite sums combine errors additively; finite products use
\eqref{eq:producterror}; a reciprocal uses its supplied positive
lower bound. The human proof contains finitely many such kinds of
steps. Their composition is a PR expression in the original
numerical parameters. This is the ordinary finite-error
translation of the proof, and it uses exactly the invariants above.

Consequently the calculation proving~\eqref{eq:Checkdefinition}
is a proof by quantifier-free induction with \(c,B,k\) as free
parameters. Universal introduction gives~\eqref{eq:checktarget}.
This establishes the claimed human proof in \(\ipra\).
Writing the corresponding deeply embedded derivation in Lean is
a separate, substantial implementation task described next.

\paragraph{Two distinctions that matter.}
PR computability alone does not establish correctness in iPRA:
the finite invariants and their proofs are required, and have been
identified above. Likewise a PR family of proofs of separate numeral
instances does not supply a proof of their universal closure.
The argument here is parameter-uniform and does not invoke an
internal reflection principle.

\subsection{Specialization to fixed primitive recursive names}

For a fixed PR complex name \(a(n)\) with a supplied regularity
hypothesis and bound \(B\), evaluate the finite routines on
a sufficiently accurate rational \(a(m)\), using an enlarged
box \(A=B+3\). Here is an explicit index rule.
Keep every Taylor degree and cutoff determined by the common
box and precision, and use exact polynomial sums with no
optional rounding. Each resulting approximant, as a function
of \(s\), is a rational complex polynomial
\(P(s)=\sum_j a_j s^j\): the rational \(\pi\) and logarithm
samples are constants independent of \(s\).
Compute, by a finite coefficient fold,
\[
 K=1+\left\lceil
       \sum_{j\ge1}j\,\normone{a_j}(2A+2)^{j-1}
       \right\rceil.
\]
The finite difference-of-powers identity proves
\(\normone{P(s)-P(t)}\le K\normone{s-t}\)
on the enlarged box. At requested precision \(k\), first
choose the expression's uniform accuracy as \(k+4\), and take
\[
                        m=k+\lceilog(K)+8.
\]
The regularity hypothesis and this polynomial bound absorb
the input error; the uniform tails for \(J\) and \(\gamrec\)
absorb the finite-expression error. This gives the indicated
named value with the requested accuracy.

This procedure is primitive recursive for each fixed input
definition. If ordinary pointwise division is required, include
an apartness witness, such as \(|s-1|\ge2^{-a}\); its inverse
accuracy is then explicit. The everywhere-defined
identity~\eqref{eq:entireFE} requires no such witness.
The specialization is a theorem schema in the fixed PR definitions,
not a first-order quantification over arbitrary real numbers.

\section{Lean formalization: a concrete completion path}\label{sec:Lean}

\subsection{Inspected ProveIt interfaces and provenance}

The GitHub reads were pinned to ProveIt commit
\code{c5093c20b5bffb0333fd42915b7bc858d5aa45cc}.
The root toolchain specifies Lean \code{v4.32.0}, and the Lake
configuration requests Mathlib \code{v4.32.0}.
The relevant namespace is \code{LeanProofs.ConciseIPRA}.
The directly inspected files and useful interfaces are
\cite{ProveItPR,ProveItBridge,ProveItReadme}:

\begin{center}
\begin{tabular}{p{0.36\textwidth}p{0.54\textwidth}}
\toprule
Source & Relevant interface\\
\midrule
\code{PR.lean} &
\code{PR}, \code{PTm}, \code{PRAAx}, and
\code{PRA.qfInd}; the induction axiom requires a QF matrix.\\
\code{BiInterpretation.lean} &
\code{biInterp}, \code{teq_iff_pra}, and
\code{pra_iff_teq}.\\
\code{NativeBiInterpretation.lean} &
\code{NFm.ofPRA}, \code{NFm.toPRA},
\code{pra_iff_native}, and \code{native_iff_pra}.\\
\bottomrule
\end{tabular}
\end{center}

The package README reports builds without \code{sorry} or added
axioms. That is a source-reported status; it was not independently
rebuilt in this environment. The metatheoretic Lean dependency of
the native bridge includes \code{Classical.choice}, while the
ordinary presentation equivalence is reported to use
\code{propext} and \code{Quot.sound}. Such ambient metatheory does
not add excluded middle to the deeply embedded object theory.

The current \code{docs/drafts} directory was inventoried. Its
analysis archives are not imported as analytic axioms in this
proof; the report proves the needed theta, gamma, and Mellin facts
directly. The archive includes a source manifest with pinned URLs
and hashes for the specific files inspected.

\subsection{Module order and exact acceptance criteria}

\begin{longtable}{p{0.22\textwidth}p{0.68\textwidth}}
\toprule
Proposed module & Work to complete and acceptance criterion\\
\midrule
\endhead
\code{Coding} &
Define fixed PR codes for signed rational and complex arithmetic,
dyadics, finite sequences, sums, products, and characteristic
functions. Prove the arithmetic identities as
\code{Prf PRAAx} derivations, and separately prove semantic
correctness in ordinary Lean. A Lean \code{Rat} implementation alone
does not meet this criterion.\\
\code{Elementary} &
Compile the specified exp, log, trigonometric, reciprocal, and
rational square-root algorithms. The output must include
object-theory proofs of their finite error sentences, with
intermediate magnitude bounds.\\
\code{Quadrature} &
Formalize common-refinement estimates and finite telescoping,
integration by parts, and the two substitutions actually used.
Each analytic shorthand must be eliminated by its finite
rational estimate.\\
\code{GaussianTheta} &
Implement Fourier quadrature, moment tails, finite Fej\'er algebra,
and the \(t=1\) normalization. Prove the parameter-uniform
theta residual sentence using~\eqref{eq:Fejercutoffs}.\\
\code{ReciprocalGamma} &
Formalize finite Euler products, their \(C_B/N\) tail,
the finite sine identity, reflection, duplication, and
Euler-integral agreement. Every inversion requires a numerical
apartness witness.\\
\code{Mellin} &
Implement \(\Jh\) exactly as in~\eqref{eq:Jhat}.
Prove Proposition~\ref{prop:Jcontract}, the half-plane
Dirichlet identification, and the regularized splitting.
This module establishes that the symmetric expression is the
zeta continuation.\\
\code{FunctionalEquation} &
Compile the fixed three-argument checker, prove its universally
closed correctness sentence in the object theory, and apply
\code{pra_iff_native} to transfer the exact result to the
native calculus.\\
\code{Semantics} &
Optionally identify the constructed names with Mathlib's analytic
objects. Keep this ordinary semantic validation separate from
the object-theory proof.\\
\bottomrule
\end{longtable}

The definitive endpoint is an actual term of
\begin{lstlisting}
Prf PRAAx [] (residualSentence compiledCheck)
\end{lstlisting}
where \code{compiledCheck : PR 3} is the fixed implementation
specified in Section~\ref{sec:arithmetic}.
The existing bridge then yields
\begin{lstlisting}
NPrf [] (NFm.ofPRA (residualSentence compiledCheck)).
\end{lstlisting}
An ambient Lean theorem saying that an ordinary function has
the correct numerical result does not replace this derivation.

\subsection{Included Lean drafts and their status}

\code{lean/AlgebraicCore.lean} records the polynomial symmetry:
\begin{lstlisting}
def xiNumerator (J : K -> K) (s : K) : K :=
  1 + s * (s - 1) * (J s + J (1 - s))

theorem xiNumerator_one_sub (J : K -> K) (s : K) :
    xiNumerator J (1 - s) = xiNumerator J s := by
  have h : (1 : K) - (1 - s) = s := by ring
  simp only [xiNumerator, h]
  ring
\end{lstlisting}
This is the easy algebraic part over a commutative ring.
It does not assert that \(J\) is an integral or that the analytic
theorem has been formalized.

\code{lean/ArithmeticBridge.lean} defines the quantifier-free
matrix \(\mathsf{check}=1\), its threefold universal closure,
and a theorem transferring a supplied iPRA derivation through
\code{pra_iff_native}.
The construction of \code{compiledCheck} and its correctness
derivation remains to be implemented.

Both files are labelled uncompiled interface drafts. No Lean
executable was available in this environment, so no kernel-checking
claim is made for them. They contain no added \code{axiom} or
\code{sorry}, but this syntactic fact does not amount to completion
of the analytic proof.

\subsection{Mathlib as a semantic cross-check}

The inspected official Mathlib source has the following theorems
\cite{MathlibZeta}:
\begin{itemize}
\item \code{completedRiemannZeta_one_sub};
\item \code{completedRiemannZeta}\textsubscript{0}\code{_one_sub}
      (the identifier uses the Unicode subscript zero);
\item \code{riemannZeta_one_sub};
\item \code{zeta_eq_tsum_one_div_nat_add_one_cpow}.
\end{itemize}
These are useful comparison theorems in ambient \(\C\).
Their noncomputable complex-analytic development is not an iPRA
proof. Mathlib also assigns totalized values at singular points;
the denominator-cleared statement avoids confusing those values
with the meromorphic assertion.

\section{Validation and delivered artifacts}\label{sec:validation}

The proof and bounds were independently reviewed in two directions:
the theta argument was checked against the gamma construction, and
the gamma argument against the theta construction. The review
covered constants, endpoint errors, zero-factor handling, and the
quantifier-free-induction boundary. The resulting argument has
been incorporated into this report.

\subsection{Exact finite checks}

The executable \code{code/quantitative_bounds.py} uses exact
Python \code{Fraction} arithmetic to check the displayed budget
inequalities at 104 selected \(J\) inputs and 91 selected theta
inputs, and a sample Fourier quadrature input.
All passed. It also reports the enormous sample workload without
attempting the full grid.
The output is in \code{quantitative_bounds_check.json}.
These checks validate selected arithmetic instances and the
transcription of the formulas; the universal proof is the
symbolic argument in the report.

\subsection{Numerical and finite-algebra diagnostics}

At 80 decimal digits, the separate executable
\code{code/numerical_diagnostics.py} compares the constructed
theta-integral numerator with \((s-1)\zeta(s)\), checks the
entire functional equation and theta transformation, and
compares two evaluations of \(J\).
The sample points include \(0,1,2,-2,-4\) and seven other
real or complex inputs.
At \(s=1\), it compares with the numerator value \(1\);
it never computes \(0\cdot\infty\).

\begin{center}
\begin{tabular}{lr}
\toprule
Diagnostic & Maximum observed discrepancy\\
\midrule
\(Z(s)\) versus \((s-1)\zeta(s)\) & \(1.624\times10^{-72}\)\\
Entire equation, absolute residual & \(2.709\times10^{-57}\)\\
Entire equation, normalized residual & \(2.166\times10^{-80}\)\\
Theta transformation & \(2.109\times10^{-81}\)\\
Two evaluations of \(J\) & \(2.330\times10^{-83}\)\\
Finite reflection identity & \(3.400\times10^{-80}\)\\
Finite duplication identity & \(5.216\times10^{-80}\)\\
\bottomrule
\end{tabular}
\end{center}

The larger absolute residual reflects large factors at one of the
complex sample points; its normalized residual is at the numerical
working precision. Exact SymPy polynomial checks of the finite
recurrence, reflection, and duplication identities also passed
for \(N=1,\ldots,8\).

The script and saved JSON expressly identify numerical tests as
diagnostics. They use floating-point special-function routines and
are not certified enclosures, iPRA proof objects, or substitutes
for the estimates proved above.

\subsection{Contents and reproduction}

The archive contains this \LaTeX\ source and its compiled PDF,
the two Lean interface drafts, the detailed formalization plan and
finite arithmetic contract, runnable diagnostic scripts and saved
outputs, a source manifest, and a README with reproduction commands.
The main report compiles with a conventional \code{pdflatex}
installation. The arithmetic budget script uses the Python standard
library; the numerical diagnostic additionally uses \code{mpmath}
and \code{sympy}.

The proof supplies the requested human-readable iPRA argument under
the explicit first-order interpretation, together with the concrete
finite contracts needed for a full formalization. The remaining
Lean work is to encode and verify those contracts inside the existing
ConciseIPRA proof system, following the module order above.

\begingroup
\small\raggedright
\begin{thebibliography}{99}

\bibitem{DLMF25}
NIST Digital Library of Mathematical Functions,
\emph{Section 25.4: Reflection Formulas}.
\url{https://dlmf.nist.gov/25.4}.
Normalizations of the zeta functional equation and the \(\xi\)-function.

\bibitem{DLMF5}
NIST Digital Library of Mathematical Functions,
\emph{Section 5.5: Functional Relations}.
\url{https://dlmf.nist.gov/5.5}.
Gamma recurrence, reflection, and duplication.

\bibitem{DLMF58}
NIST Digital Library of Mathematical Functions,
\emph{Section 5.8: Infinite Products}.
\url{https://dlmf.nist.gov/5.8}.
Euler's limit and reciprocal-gamma product.

\bibitem{DLMF422}
NIST Digital Library of Mathematical Functions,
\emph{Section 4.22: Infinite Products and Partial Fractions}.
\url{https://dlmf.nist.gov/4.22}.
Euler's sine product.

\bibitem{ProveItReadme}
Vladimir Reshetnikov, \emph{ProveIt: ConciseIPRA README},
snapshot \code{c5093c20b5bffb0333fd42915b7bc858d5aa45cc}.
\url{https://github.com/VladimirReshetnikov/ProveIt/blob/c5093c20b5bffb0333fd42915b7bc858d5aa45cc/Logic/PrimitiveRecursiveArithmetic/ConciseIPRA/README.md}.

\bibitem{ProveItPR}
Vladimir Reshetnikov, \emph{ProveIt: ConciseIPRA/PR.lean},
same snapshot.
\url{https://github.com/VladimirReshetnikov/ProveIt/blob/c5093c20b5bffb0333fd42915b7bc858d5aa45cc/Logic/PrimitiveRecursiveArithmetic/ConciseIPRA/Lean/ConciseIPRA/PR.lean}.

\bibitem{ProveItBridge}
Vladimir Reshetnikov,
\emph{ProveIt: ConciseIPRA/NativeBiInterpretation.lean},
same snapshot.
\url{https://github.com/VladimirReshetnikov/ProveIt/blob/c5093c20b5bffb0333fd42915b7bc858d5aa45cc/Logic/PrimitiveRecursiveArithmetic/ConciseIPRA/Lean/ConciseIPRA/NativeBiInterpretation.lean}.

\bibitem{MathlibZeta}
The Mathlib community, \emph{Mathlib/NumberTheory/LSeries/RiemannZeta.lean},
tag \code{v4.32.0}.
\url{https://github.com/leanprover-community/mathlib4/blob/v4.32.0/Mathlib/NumberTheory/LSeries/RiemannZeta.lean}.
See also the official generated documentation:
\url{https://leanprover-community.github.io/mathlib4_docs/Mathlib/NumberTheory/LSeries/RiemannZeta.html}.

\end{thebibliography}
\endgroup
\end{document}

